% ============================================================
%  GelSight 光度立体重建算法精度问题分析报告
%  编译命令：xelatex -> bibtex -> xelatex -> xelatex
%  依赖：TeX Live 2021+（含 ctex、listings、tcolorbox 等包）
% ============================================================
\documentclass[12pt, a4paper]{ctexart}

% ---- 页面尺寸 ----
\usepackage[top=2.8cm, bottom=2.8cm, left=3.0cm, right=3.0cm]{geometry}

% ---- 行距 ----
\usepackage{setspace}
\onehalfspacing

% ---- 数学 ----
\usepackage{amsmath, amssymb, amsthm}

% ---- 代码块 ----
\usepackage{listings}
\usepackage{xcolor}
\definecolor{codebg}{RGB}{248,248,248}
\definecolor{codeframe}{RGB}{200,200,200}
\definecolor{keywordcolor}{RGB}{0,0,180}
\definecolor{commentcolor}{RGB}{100,100,100}
\definecolor{stringcolor}{RGB}{0,130,0}
\lstset{
  language        = Python,
  basicstyle      = \ttfamily\small,
  keywordstyle    = \color{keywordcolor}\bfseries,
  commentstyle    = \color{commentcolor}\itshape,
  stringstyle     = \color{stringcolor},
  showstringspaces= false,
  breaklines      = true,
  frame           = single,
  framesep        = 5pt,
  rulecolor       = \color{codeframe},
  backgroundcolor = \color{codebg},
  xleftmargin     = 1.5em,
  xrightmargin    = 0.5em,
  numbers         = left,
  numberstyle     = \tiny\color{gray},
  stepnumber      = 1,
  numbersep       = 6pt,
  captionpos      = b,
  aboveskip       = 0.8em,
  belowskip       = 0.5em,
}

% ---- 表格 ----
\usepackage{booktabs}
\usepackage{tabularx}
\usepackage{multirow}

% ---- 超链接 ----
\usepackage{hyperref}
\hypersetup{
  colorlinks = true,
  linkcolor  = {blue!65!black},
  citecolor  = {blue!65!black},
  urlcolor   = {blue!65!black},
  pdftitle   = {GelSight 光度立体重建算法精度问题分析},
}

% ---- 参考文献 ----
\usepackage[numbers, sort&compress]{natbib}

% ---- 彩色提示框 ----
\usepackage{tcolorbox}
\tcbuselibrary{skins, breakable}

\newtcolorbox{bugbox}[1][]{
  colback   = red!4!white,
  colframe  = red!60!black,
  fonttitle = \bfseries\small,
  title     = {Bug / 精度缺陷},
  #1,
  left=5pt, right=5pt, top=4pt, bottom=4pt,
  boxrule   = 0.7pt,
  breakable,
}
\newtcolorbox{fixbox}[1][]{
  colback   = green!4!white,
  colframe  = green!50!black,
  fonttitle = \bfseries\small,
  title     = {修复方案},
  #1,
  left=5pt, right=5pt, top=4pt, bottom=4pt,
  boxrule   = 0.7pt,
  breakable,
}

% ---- 页眉/页脚 ----
\usepackage{fancyhdr}
\pagestyle{fancy}
\fancyhf{}
\rhead{\small\textit{GelSight 光度立体算法精度分析}}
\lhead{\small\textit{内部技术报告}}
\cfoot{\thepage}

% ---- 杂项 ----
\usepackage{enumitem}
\setlist[enumerate]{leftmargin=*, label=\arabic*.}
\setlist[itemize]{leftmargin=*}

% ============================================================
\begin{document}

% ---------- 封面 ----------
\begin{titlepage}
  \centering
  \vspace*{2cm}
  {\LARGE\bfseries GelSight 光度立体重建算法\\[0.4em]精度问题分析报告\par}
  \vspace{0.6em}
  {\large\itshape Accuracy Issue Analysis of the GelSight\\
   Photometric Stereo Reconstruction Pipeline\par}
  \vspace{2em}
  \rule{0.6\linewidth}{0.4pt}\\[1em]
  {\normalsize 基于源码 \texttt{gelsight\_ps} 的逐模块审查\\[0.3em]
   覆盖模块：\texttt{geometry.py}、\texttt{ball\_calibrate.py}、\\
   \texttt{poisson.py}、\texttt{reconstruct.py}、\texttt{image\_loader.py}\par}
  \vspace{2em}
  {\normalsize\today}
  \vfill
  \begin{tcolorbox}[colback=gray!8, colframe=gray!50, width=0.85\linewidth,
                    boxrule=0.5pt, left=8pt, right=8pt]
    \small\textbf{摘要：}本报告对 \texttt{gelsight\_ps} 软件包的光度立体标定与重建
    流程进行全面精度审查，识别出 \textbf{6 类、12 个具体问题}，包括 3 个可立即修
    复的代码逻辑缺陷及 9 个需算法层面改进的问题。各问题均附有定位代码、
    误差分析、相关文献及修复建议，并按影响程度提供修复优先级排序。
  \end{tcolorbox}
\end{titlepage}

% ---------- 目录 ----------
\tableofcontents
\newpage

% ============================================================
\section{引言}
\label{sec:intro}
% ============================================================

GelSight 系列触觉传感器~\cite{johnson2009retrographic,yuan2017gelsight} 通过观察弹性
凝胶层在接触时的光度变化，配合光度立体（Photometric Stereo, PS）
算法~\cite{woodham1980photometric} 重建接触面的三维形貌。其标定流程的核心是：
用已知几何形状（钢球）压入传感器，建立 RGB 响应到表面法向的逐像素映射（LUT）；
推理时通过 LUT 预测法向，再经 Frankot-Chellappa 频域积分~\cite{frankot1988method}
恢复深度图。

\texttt{gelsight\_ps} 软件包实现了上述完整流程。经逐行源码审查，当前实现在以下
六个方向存在精度问题，本报告逐一论述，并结合近年文献给出改进方向。

\begin{table}[h]
  \centering
  \caption{问题分类概览}
  \label{tab:overview}
  \begin{tabular}{clc}
    \toprule
    \textbf{章节} & \textbf{问题类别} & \textbf{问题数} \\
    \midrule
    \S\ref{sec:gt_bias}    & 地面真值法向量系统性偏差   & 2 \\
    \S\ref{sec:regression} & 回归模型内在局限           & 4 \\
    \S\ref{sec:integration}& 法向积分误差               & 2 \\
    \S\ref{sec:bugs}       & 代码逻辑缺陷（可立即修复） & 2 \\
    \S\ref{sec:eval}       & 评估指标非标准             & 1 \\
    \S\ref{sec:preproc}    & 预处理缺陷                 & 1 \\
    \midrule
    & \textbf{合计} & \textbf{12} \\
    \bottomrule
  \end{tabular}
\end{table}

% ============================================================
\section{地面真值法向量的系统性偏差}
\label{sec:gt_bias}
% ============================================================

\subsection{\texttt{nz\_floor} 直接篡改训练标签}
\label{subsec:nz_floor}

\subsubsection*{问题定位}

\texttt{geometry.py}，函数 \texttt{sphere\_normals\_for\_contact()}，第 59--66 行：

\begin{lstlisting}[caption={\texttt{nz\_floor} 对训练标签的错误应用}]
nz_raw = z[inside] / R

# 本意是防止推理阶段的梯度奇异，但错误地修改了训练标签
nz[inside] = np.maximum(nz_raw, float(nz_floor))  # nz_floor = 0.12

n = np.stack([nx, ny, nz], axis=-1)
norm = np.linalg.norm(n, axis=-1, keepdims=True) + 1e-9
n = n / norm                                       # 归一化后 nz 不再是球面真值
\end{lstlisting}

\begin{bugbox}
当接触区边缘像素的 $n_z^{raw} < 0.12$ 时，代码将其强制提升至 $0.12$，随后对向量
$(n_x,\, n_y,\, 0.12)$ 做归一化，得到：
\[
  \hat{\mathbf{n}}_{label}
  = \frac{(n_x,\; n_y,\; 0.12)}{\|(n_x,\; n_y,\; 0.12)\|}
  \neq \hat{\mathbf{n}}_{true}
\]
回归器以 $\hat{\mathbf{n}}_{label}$ 为目标训练，导致高倾斜角区域
（$\theta > \arccos(0.12) \approx 83^\circ$）的法向预测出现系统性偏差。
\end{bugbox}

\subsubsection*{理论依据}

Woodham~\cite{woodham1980photometric} 在原始 PS 论文中即指出：近掠射角样本的法向
不确定度高，应\textbf{从训练集中剔除}而非修正。正确的处理策略是扩大
\texttt{rim\_shrink\_pix}，使进入训练的像素均满足 $n_z^{raw} \geq \varepsilon$，
但不对法向向量本身做任何修改。

\begin{fixbox}
\texttt{nz\_floor} 作为梯度计算时的分母截断（防止 $p = n_x/n_z$ 奇异化），
应仅在\textbf{推理阶段} \texttt{reconstruct.py} 中使用；标定时只剔除，不修正：

\begin{lstlisting}
# geometry.py -- 修复方案：排除低 nz 像素，而非修改其标签
valid_nz = (z[inside] / R) >= float(nz_floor)
inside_valid = inside.copy()
inside_valid[inside] = valid_nz        # 只保留 nz_raw >= nz_floor 的像素

mask = inside_valid.astype(np.uint8)

# nz 直接使用原始球面法向，不做任何 floor 截断
nz[inside_valid] = z[inside][valid_nz] / R
\end{lstlisting}
\end{fixbox}

% ----------------------------------------
\subsection{Hertz 接触变形导致球面几何模型失效}
\label{subsec:hertz}

\subsubsection*{问题描述}

\texttt{sphere\_normals\_for\_contact()} 假设接触区表面法向严格等于刚性球面的法向。
然而，GelSight 所用硅胶（通常为 Shore~00-20 至 00-40 级）是弹性软材料。
根据 Hertz 弹性接触理论，球体压入弹性半空间时，接触边缘的凝胶会发生非线性
弯折变形，真实法向与刚性球面法向之间可差 $5°$--$15°$，
Yuan~et~al.~\cite{yuan2017gelsight} 及 Dong~et~al.~\cite{dong2017improved}
对此有定性讨论。

\subsubsection*{改进方向}

\begin{enumerate}
  \item \textbf{多形状交叉验证}：使用平面、圆柱、球面等不同已知曲率的标准件
        进行标定数据采集，当各形状的预测深度图与几何真值一致时，
        说明标定模型已收敛至物理正确的状态。
  \item \textbf{数据驱动的弹性修正因子}：在 $\hat{\mathbf{n}}_{sphere}$ 基础上
        引入与接触深度相关的法向修正量，可通过对已知平面的残差法向场进行
        统计建模得到。
\end{enumerate}

% ============================================================
\section{回归模型的内在局限}
\label{sec:regression}
% ============================================================

\subsection{线性（二阶多项式）模型无法描述复杂光学效应}
\label{subsec:linear_model}

\subsubsection*{当前特征空间}

\texttt{build\_feature\_extractor()} 构造 11 维多项式特征向量，对应线性回归：
\[
  \boldsymbol{\phi}(\mathbf{c}) = \bigl[
    r',\; g',\; b',\; I,\; 1,\; R^2,\; G^2,\; B^2,\; RG,\; RB,\; GB
  \bigr], \qquad
  \hat{\mathbf{n}} = \mathbf{W}\boldsymbol{\phi}
\]
其中 $r' = R/(R+G+B)$，$I = R+G+B$。

\subsubsection*{非线性效应清单}

\begin{table}[h]
  \centering
  \caption{GelSight 光学通道中的非线性效应及线性模型的描述能力}
  \label{tab:nonlinear}
  \begin{tabular}{lcc}
    \toprule
    \textbf{效应}                      & \textbf{线性可描述？} & \textbf{影响区域} \\
    \midrule
    多 LED 距离衰减（$1/d^2$ 律）      & 否                    & 全局              \\
    凝胶--空气界面菲涅尔反射            & 否                    & 大倾斜角区域      \\
    硅胶内次表面散射                    & 否                    & 接触核心区        \\
    多 LED 混合色散差异                 & 部分                  & 全局              \\
    \bottomrule
  \end{tabular}
\end{table}

\subsubsection*{文献支撑与改进方向}

Lambeta~et~al.~\cite{lambeta2020digit} 在 DIGIT 传感器中用卷积神经网络替代
逐像素线性回归，在高曲率区域的法向误差显著降低。
短期内可行的改进是扩展特征空间（加入三阶项）：

\begin{lstlisting}[caption={短期改进：在特征中加入三阶单项式}]
feature_names += ["R^3", "G^3", "B^3"]

# make_features_from_rgb 中对应追加
feats += [R * R * R, G * G * G, B * B * B]
\end{lstlisting}

% ----------------------------------------
\subsection{8×8 分块网格过粗，无法捕捉像素级光照变化}
\label{subsec:block_grid}

\subsubsection*{问题定位}

\texttt{ball\_calibrate.py}，第 1066 行：

\begin{lstlisting}[caption={硬编码的 8×8 分块}]
K  = 8
bh = (H + K - 1) // K   # 对 480x640 图像约为 60x80 像素/块
bw = (W + K - 1) // K
\end{lstlisting}

\subsubsection*{分析}

对 $480 \times 640$ 传感器，每块约 $60 \times 80$ 像素。同一块内，像素到 LED
的方向角最大可相差 $20°$--$30°$，块内真正的 RGB→法向映射不是同一个线性变换。
块级残差 $\Delta\mathbf{W}$ 仅捕捉块间的低频空间变化，块内高频变化被忽略。

Wang~et~al.~\cite{wang2021gelsightwedge} 的 GelSight~Wedge 使用逐像素 LUT
而非分块近似，这正是当前分块设计所缺失的精度层级。

\begin{fixbox}
将 $K$ 提高至 16 或 32；对覆盖度满足阈值的像素直接做逐像素岭回归，
覆盖不足的像素用 RBF 插值填充：

\begin{lstlisting}
K = 16   # 或 32，块大小降至约 15x20 像素

# 进一步：对 Ns >= min_spp 的像素直接逐像素回归
for y, x in zip(*np.where(Ns >= min_spp)):
    W_field[y, x] = np.linalg.solve(
        XtX[y, x] + lam * np.eye(FEAT_DIM),
        XtY[y, x]
    )
# 其余像素用邻域加权（RBF）插值填充
\end{lstlisting}
\end{fixbox}

% ----------------------------------------
\subsection{高斯融合邻域搜索范围不足}
\label{subsec:gaussian_blend}

\subsubsection*{问题定位}

\texttt{ball\_calibrate.py}，第 1113--1125 行：

\begin{lstlisting}[caption={sigma 与搜索范围不匹配}]
sigma_x = bw * 0.75      # 有效半径约 1.5 个块宽
sigma_y = bh * 0.75

# 但搜索范围仅 ±1 块（约 ±1.33 sigma），应至少搜索 ±2 sigma
for by in range(max(0, int(by_f) - 1), min(K, int(by_f) + 2)):
    for bx in range(max(0, int(bx_f) - 1), min(K, int(bx_f) + 2)):
\end{lstlisting}

\subsubsection*{后果}

图像角点处（如第 $0$ 行第 $0$ 列）只有 1--2 个有效邻居，$w_{sum}$ 极小，导致：
\[
  \mathbf{W}_{acc}
  = \mathbf{W}_{global} + \frac{\mathbf{W}_{acc} - \mathbf{W}_{global}}{w_{sum}}
  \xrightarrow{w_{sum} \to 0} \text{残差被无限放大}
\]

\begin{fixbox}
搜索范围扩展至 $\lceil 2\sigma / \text{block\_size} \rceil = 2$ 块，
同时增加 $w_{sum} < \varepsilon$ 的保护性退化：

\begin{lstlisting}
radius = 2   # 扩展为 ±2 块
for by in range(max(0, int(by_f) - radius),
                min(K, int(by_f) + radius + 1)):
    for bx in range(max(0, int(bx_f) - radius),
                    min(K, int(bx_f) + radius + 1)):
        ...

# 归一化保护
if wsum > 1e-6:
    W_field[y, x] = W_global + (W_acc - W_global) / wsum
else:
    W_field[y, x] = W_global   # 无有效邻居时退化为全局模型
\end{lstlisting}
\end{fixbox}

% ----------------------------------------
\subsection{样本积累的像素级 Python 循环效率极低}
\label{subsec:loop}

\texttt{ball\_calibrate.py}，第 1009--1026 行通过 Python 循环逐像素累积
$\mathbf{X}^T\mathbf{X}$ 和 $\mathbf{X}^T\mathbf{Y}$，每帧接触区约
3\,000--10\,000 像素、$F=11$ 维特征，全部在 Python 层迭代。
此问题\textbf{不影响精度}，但标定速度可提升 10--50 倍：

\begin{lstlisting}[caption={向量化方案（取代像素级 for 循环）}]
ys, xs = np.where(core > 0)
v  = X[ys, xs]              # [N, F]
yv = Y[ys, xs]              # [N, out_dim]
w  = weights[ys, xs]        # [N]

# XtX 批量更新
np.add.at(XtX, (ys, xs),
          np.einsum('ni,nj->nij', v * w[:, None], v))

# XtY 批量更新
np.add.at(XtY, (ys, xs),
          np.einsum('ni,nj->nij', v * w[:, None], yv))

Ns[ys, xs] += (w > 0).astype(np.int32)
\end{lstlisting}

% ============================================================
\section{法向积分误差}
\label{sec:integration}
% ============================================================

\subsection{Frankot-Chellappa 的周期性假设引发低频振铃}
\label{subsec:fc_ringing}

\subsubsection*{算法局限}

\texttt{poisson.py} 实现的 Frankot-Chellappa 积分~\cite{frankot1988method}
通过 2D FFT 求解 Poisson 方程：
\[
  Z(\omega_x, \omega_y)
  = \frac{-j\omega_x\, P(\omega_x, \omega_y)
          -j\omega_y\, Q(\omega_x, \omega_y)}
         {\omega_x^2 + \omega_y^2}
\]

FFT 隐含地假设信号在边界处首尾周期连接。深度图的左右上下边界值通常不相等，
这一``跳跃''引发 Gibbs 效应（低频振铃伪影），在平坦背景区最为明显。

\subsubsection*{文献支撑}

Qu\'eau~et~al.~\cite{queeau2018normal} 的综述系统梳理了这一问题，
推荐基于离散余弦变换（DCT-II）的 Poisson 求解器~\cite{simchony1990direct}，
其满足 Neumann 零通量边界条件，从根本上消除振铃：
\[
  \nabla^2 z = \nabla \cdot \mathbf{v}, \qquad
  \left.\frac{\partial z}{\partial \mathbf{n}}\right|_{\partial\Omega} = 0
  \quad (\text{Neumann BC})
\]

Harker \& O'Leary~\cite{harker2008least} 进一步提出最小二乘法，
对不可积梯度场给出最优（最小 Frobenius 范数误差）积分结果。

\begin{fixbox}[title={修复方案：基于 DCT-II 的 Poisson 求解器}]
\begin{lstlisting}
from scipy.fft import dctn, idctn

def poisson_dct(p: np.ndarray, q: np.ndarray) -> np.ndarray:
    """Neumann-BC Poisson integration via DCT-II.

    Replaces frankot_chellappa() in poisson.py.
    Reference: Simchony et al., IEEE TPAMI 1990.
    """
    H, W = p.shape
    # 散度 div = dp/dx + dq/dy（前向有限差分）
    div = np.zeros((H, W), dtype=np.float64)
    div[:, 1:]  += p[:, :-1]
    div[:, :-1] -= p[:, :-1]
    div[1:, :]  += q[:-1, :]
    div[:-1, :] -= q[:-1, :]
    # DCT-II 域 Poisson 求解
    D    = dctn(div, norm='ortho')
    cy   = np.cos(np.pi * np.arange(H) / H)[:, None]
    cx   = np.cos(np.pi * np.arange(W) / W)[None, :]
    denom = 2.0 * (cy - 1.0) + 2.0 * (cx - 1.0)
    denom[0, 0] = 1.0
    D /= denom
    D[0, 0] = 0.0    # DC 分量（深度基线任意）
    return idctn(D, norm='ortho').astype(np.float32)
\end{lstlisting}
\end{fixbox}

% ----------------------------------------
\subsection{低 $n_z$ 处梯度放大污染全局积分}
\label{subsec:nz_amplification}

\subsubsection*{问题定位}

\texttt{reconstruct.py}，推理路径：

\begin{lstlisting}[caption={nz=0.1 时梯度被放大 10 倍}]
p_grad = n[..., 0] / np.maximum(n[..., 2], 0.1)
q_grad = n[..., 1] / np.maximum(n[..., 2], 0.1)
\end{lstlisting}

当 $n_z = 0.1$（约 $84°$ 倾角）时，梯度值 $p = n_x/0.1$ 被放大 10 倍。
由于 FFT 是全局算子，\textbf{少量噪声大的边缘梯度会将误差扩散至整幅深度图}，
使接触区以外的背景也出现波纹。

\subsubsection*{文献支撑}

Harker \& O'Leary~\cite{harker2008least} 及 Qu\'eau~et~al.~\cite{queeau2018normal}
均推荐使用置信度权重对不可靠像素降权，置信度常取 $c = n_z^2$：

\begin{fixbox}
\begin{lstlisting}
# 置信度权重：nz 越低，权重越小
nz_conf = np.clip(n[..., 2], 0.0, 1.0)
conf    = nz_conf ** 2                  # 可调整指数

# 更保守的截断阈值，同时乘置信权重
nz_safe = np.maximum(n[..., 2], 0.25)  # 截断 >= 0.25 (约 76 度)
p_grad  = (n[..., 0] / nz_safe) * conf
q_grad  = (n[..., 1] / nz_safe) * conf
\end{lstlisting}
\end{fixbox}

% ============================================================
\section{代码逻辑缺陷}
\label{sec:bugs}
% ============================================================

\subsection{\texttt{image\_dir} 模式下掩膜被施加两次}
\label{subsec:double_mask}

\subsubsection*{问题定位}

\texttt{reconstruct.py}，\texttt{run\_reconstruction()} 的 \texttt{image\_dir} 分支：

\begin{lstlisting}[caption={接触掩膜被连续施加两次，导致 d\_diff 计算错误}]
# 第一次：对绝对深度施加掩膜
if contact_mask_enable and ref_lin is not None:
    mask = _residual_mask(...)
    if mask is not None:
        depth = depth * (mask / 255.0)        # depth 已被掩膜

# 之后：
d_diff = depth - ref_depth                    # depth 已不完整（非 depth_orig）

# 第二次：对 d_diff 再次施加掩膜
if contact_mask_enable and ref_lin is not None:
    mask = _residual_mask(...)
    if mask is not None:
        d_diff = d_diff * (mask / 255.0)      # 等效二次掩膜
\end{lstlisting}

\begin{bugbox}
实际计算结果为：
\[
  d\_diff = (depth_{orig} \cdot M - ref\_depth) \cdot M
\]
而正确结果应为：
\[
  d\_diff = (depth_{orig} - ref\_depth) \cdot M
\]

在掩膜外缘，第一次掩膜将 $depth_{orig}$ 置零，使得
$(0 - ref\_depth) \cdot M = -ref\_depth \cdot M$；
若 $ref\_depth \neq 0$，掩膜边界处会引入假深度差。
\end{bugbox}

\begin{fixbox}
移除对绝对 \texttt{depth} 的第一次掩膜，仅对 \texttt{d\_diff} 施加一次：

\begin{lstlisting}
# 先计算相对深度（均使用未掩膜的原始深度）
d_diff = (depth - ref_depth).astype(np.float32)

# 再施加一次掩膜
if contact_mask_enable and ref_lin is not None:
    mask = _residual_mask(bgr_lin, ref_lin,
                          cm_w_chroma, cm_w_int,
                          cm_thr_rel, cm_min_area)
    if mask is not None:
        d_diff *= (mask.astype(np.float32) / 255.0)
\end{lstlisting}
\end{fixbox}

% ----------------------------------------
\subsection{\texttt{out\_dim=2} 路径的 \texttt{sqrt(负数)} $\to$ 静默 \texttt{NaN}}
\label{subsec:sqrt_nan}

\subsubsection*{问题定位}

\texttt{reconstruct.py}，\texttt{\_lut\_predict\_normals()}：

\begin{lstlisting}[caption={clip 在 sqrt 之后无法清除 NaN}]
nz = np.sqrt(1.0 - nx * nx - ny * ny)  # 当 nx^2 + ny^2 > 1 时 -> NaN
nz = np.clip(nz, 0.0, 1.0)             # clip(NaN, 0, 1) 仍然是 NaN ！
\end{lstlisting}

\begin{bugbox}
当回归器在训练数据稀疏区域外推，预测 $n_x^2 + n_y^2 > 1$ 时，
\texttt{np.sqrt} 对负数返回 \texttt{NaN}，而
\textbf{\texttt{np.clip} 不清除 \texttt{NaN}}。
\texttt{NaN} 沿 normalize $\to$ $p$ grad $\to$ $q$ grad $\to$ FFT
全链路传播，最终深度图出现随机 \texttt{NaN} 块。

注：\texttt{ball\_calibrate.py} 第 1258 行的评估路径已正确处理此问题，
但 \texttt{reconstruct.py} 未同步。
\end{bugbox}

\begin{fixbox}
将 \texttt{clip} 移入 \texttt{sqrt} 内部（与 \texttt{ball\_calibrate.py} 保持一致）：

\begin{lstlisting}
# 修复：clip 在 sqrt 之前，消除负数输入
nz = np.sqrt(np.clip(1.0 - nx * nx - ny * ny, 0.0, 1.0))
nz = np.clip(nz, 0.0, 1.0)  # 保留第二行作为防御性检查
\end{lstlisting}
\end{fixbox}

% ============================================================
\section{评估指标问题}
\label{sec:eval}
% ============================================================

\texttt{ball\_calibrate.py}，第 1275 行：

\begin{lstlisting}
acc = (1.0 + cos) * 0.5    # 自定义指标，范围 [0, 1]
\end{lstlisting}

该 $\frac{1+\cos\theta}{2}$ 指标在 PS 文献中不通用，且对大角度误差不敏感：
$90°$ 误差时指标值为 $0.5$，表面上``尚可''。

PS 领域的标准评估指标为\textbf{平均角度误差（Mean Angular Error, MAE）}，
单位为度，自 Woodham~\cite{woodham1980photometric} 后续大量文献均采用此标准：
\[
  \text{MAE} = \frac{1}{|\mathcal{M}|}
  \sum_{(u,v) \in \mathcal{M}}
  \arccos\!\Bigl(\hat{\mathbf{n}}_{pred}^{(u,v)} \cdot \hat{\mathbf{n}}_{gt}^{(u,v)}\Bigr)
\]

\begin{fixbox}
\begin{lstlisting}
err_deg    = np.degrees(np.arccos(np.clip(cos, -1.0, 1.0)))
mae        = float(err_deg[mask].mean())
median_err = float(np.median(err_deg[mask]))

# 报告分段精度（与文献可比）
for thresh in [5.0, 10.0, 15.0]:
    frac = float((err_deg[mask] < thresh).mean()) * 100.0
    print(f"  Angular error < {thresh:.0f} deg: {frac:.1f}%")
\end{lstlisting}
\end{fixbox}

% ============================================================
\section{预处理缺陷：缺少平场校正}
\label{sec:preproc}
% ============================================================

\texttt{Preprocessor.apply()} 执行暗场减法和 gamma 线性化，但未校正 LED 照明的
空间不均匀性（晕影效应，vignetting）。对 GelSight 传感器，图像中心与边缘的 LED
照射强度比通常达 2--4 倍~\cite{yuan2017gelsight,dong2017improved}。

这一空间变化无法被逐像素 LUT 完全吸收，因为 LUT 以 RGB 值为输入，
而同一 RGB 值在不同图像位置可能代表完全不同的法向——只有在平场均匀化后，
RGB 值才能与法向形成唯一对应。

\begin{fixbox}
对平坦未受力凝胶面拍摄参考图 $\mathbf{I}_{flat}$，
在标定与推理时统一做平场除法（在现有暗场校正之后追加）：

\begin{lstlisting}
# Preprocessor.__init__ 中加载平场图
if "flat_field_path" in cfg and cfg["flat_field_path"]:
    self.flat_field = imread_color(cfg["flat_field_path"])
    self.flat_field = pre_dark_white(self.flat_field)  # 同样做暗/白校正
else:
    self.flat_field = None

# Preprocessor.apply 中追加
if self.flat_field is not None:
    bgr_lin = bgr_lin / (self.flat_field + 1e-4)
    bgr_lin = np.clip(bgr_lin, 0.0, None)
\end{lstlisting}
\end{fixbox}

% ============================================================
\section{问题汇总与修复优先级}
\label{sec:summary}
% ============================================================

\begin{table}[h]
  \centering
  \caption{全部问题汇总、影响评估与修复优先级}
  \label{tab:summary}
  \small
  \begin{tabularx}{\textwidth}{clXcc}
    \toprule
    \textbf{\#} & \textbf{文件} & \textbf{问题描述}
                & \textbf{精度影响} & \textbf{修复难度} \\
    \midrule
    10 & \texttt{reconstruct.py}     & \texttt{sqrt(负数)} $\to$ 静默 NaN
       & \textbf{极高}  & 极低 \\
    9  & \texttt{reconstruct.py}     & 掩膜施加两次，\texttt{d\_diff} 错误
       & 高             & 低   \\
    1  & \texttt{geometry.py}        & \texttt{nz\_floor} 污染训练标签
       & 高             & 低   \\
    7  & \texttt{poisson.py}         & FC 周期假设引发低频振铃
       & 高             & 中   \\
    8  & \texttt{reconstruct.py}     & 低 $n_z$ 梯度放大污染全局积分
       & 高             & 中   \\
    5  & \texttt{ball\_calibrate.py} & 高斯融合邻域范围不足
       & 中             & 低   \\
    4  & \texttt{ball\_calibrate.py} & $K=8$ 分块过粗
       & 中高           & 中   \\
    11 & \texttt{ball\_calibrate.py} & 评估指标非标准 MAE
       & 可观测性差     & 低   \\
    12 & \texttt{image\_loader.py}   & 无平场（晕影）校正
       & 中             & 中   \\
    2  & \texttt{geometry.py}        & Hertz 变形未建模
       & 中             & 高   \\
    3  & \texttt{ball\_calibrate.py} & 线性多项式模型不足
       & 中高           & 高   \\
    6  & \texttt{ball\_calibrate.py} & 像素循环低效（不影响精度）
       & 无             & 低   \\
    \bottomrule
  \end{tabularx}
\end{table}

\noindent\textbf{建议修复顺序（最小改动 $\to$ 最大收益）：}

\begin{enumerate}
  \item \textbf{立即可修（不改变算法结构）：}\\
        问题 \#10（\texttt{sqrt NaN}）$\to$
        问题 \#9（双重掩膜）$\to$
        问题 \#11（评估指标）$\to$
        问题 \#6（向量化加速）

  \item \textbf{短期改进（修改参数逻辑）：}\\
        问题 \#1（\texttt{nz\_floor} 标签修复）$\to$
        问题 \#5（扩大高斯融合邻域）

  \item \textbf{中期改进（替换子模块）：}\\
        问题 \#7（换 DCT 积分器）$\to$
        问题 \#8（置信度加权梯度）$\to$
        问题 \#12（添加平场校正）

  \item \textbf{长期改进（重构算法）：}\\
        问题 \#4（增大 $K$ 至逐像素回归）$\to$
        问题 \#3（引入 MLP 非线性模型）$\to$
        问题 \#2（Hertz 弹性修正）
\end{enumerate}

% ============================================================
% 参考文献
% ============================================================
\newpage
\bibliographystyle{unsrtnat}
\bibliography{references}

\end{document}
